\section{Mitigating Hallucination through Structured Grounding}
\label{sec:gr-hallucination}

Hallucination---the generation of content that is fluent but unfaithful to any
source---is a well-documented and pervasive failure mode of neural language
generation~\cite{ji2023hallucination}. Retrieval augmentation reduces it by
conditioning generation on retrieved evidence rather than parametric memory
alone~\cite{lewis2020rag, gao2023ragsurvey}, and Graph RAG sharpens this effect by
grounding the answer in explicit entities and relations whose provenance can be
traced back through the graph~\cite{peng2024graphrag}. Because a knowledge graph
supplies structured, verifiable facts, unifying it with a language model directly
targets the tendency of parametric models to fabricate
knowledge~\cite{pan2024unifying}.

Structured grounding also makes the residual problem measurable. Fine-grained
factuality metrics decompose a generated answer into atomic facts and check each
against a knowledge source~\cite{min2023factscore}, which aligns naturally with a
graph whose triples provide exactly such atomic, checkable
units~\cite{pan2024unifying}. Grounding therefore does more than lower the
hallucination rate; it provides the reference needed to detect and quantify what
remains~\cite{ji2023hallucination, min2023factscore}.
